% Options for packages loaded elsewhere
\PassOptionsToPackage{unicode}{hyperref}
\PassOptionsToPackage{hyphens}{url}
\documentclass[
  11pt,
]{article}
\usepackage{xcolor}
\usepackage[margin=24mm]{geometry}
\usepackage{amsmath,amssymb}
\setcounter{secnumdepth}{-\maxdimen} % remove section numbering
\usepackage{iftex}
\ifPDFTeX
  \usepackage[T1]{fontenc}
  \usepackage[utf8]{inputenc}
  \usepackage{textcomp} % provide euro and other symbols
\else % if luatex or xetex
  \usepackage{unicode-math} % this also loads fontspec
  \defaultfontfeatures{Scale=MatchLowercase}
  \defaultfontfeatures[\rmfamily]{Ligatures=TeX,Scale=1}
\fi
\usepackage{lmodern}
\ifPDFTeX\else
  % xetex/luatex font selection
\fi
% Use upquote if available, for straight quotes in verbatim environments
\IfFileExists{upquote.sty}{\usepackage{upquote}}{}
\IfFileExists{microtype.sty}{% use microtype if available
  \usepackage[]{microtype}
  \UseMicrotypeSet[protrusion]{basicmath} % disable protrusion for tt fonts
}{}
\makeatletter
\@ifundefined{KOMAClassName}{% if non-KOMA class
  \IfFileExists{parskip.sty}{%
    \usepackage{parskip}
  }{% else
    \setlength{\parindent}{0pt}
    \setlength{\parskip}{6pt plus 2pt minus 1pt}}
}{% if KOMA class
  \KOMAoptions{parskip=half}}
\makeatother
\usepackage{longtable,booktabs,array}
\newcounter{none} % for unnumbered tables
\usepackage{calc} % for calculating minipage widths
% Correct order of tables after \paragraph or \subparagraph
\usepackage{etoolbox}
\makeatletter
\patchcmd\longtable{\par}{\if@noskipsec\mbox{}\fi\par}{}{}
\makeatother
% Allow footnotes in longtable head/foot
\IfFileExists{footnotehyper.sty}{\usepackage{footnotehyper}}{\usepackage{footnote}}
\makesavenoteenv{longtable}
\setlength{\emergencystretch}{3em} % prevent overfull lines
\providecommand{\tightlist}{%
  \setlength{\itemsep}{0pt}\setlength{\parskip}{0pt}}
\usepackage{bookmark}
\IfFileExists{xurl.sty}{\usepackage{xurl}}{} % add URL line breaks if available
\urlstyle{same}
\hypersetup{
  pdftitle={Certified hybrid composition for exact matrix multiplication},
  pdfauthor={Alejandro Zarzuelo Urdiales},
  hidelinks,
  pdfcreator={LaTeX via pandoc}}

\title{Certified hybrid composition for exact matrix multiplication}
\author{Alejandro Zarzuelo Urdiales\\[3pt]\small\href{https://rsihouse.ai/openmath}{Open Math}}
\date{7 October 2026}

\begin{document}
\maketitle

\subsection{Abstract}\label{abstract}

We describe an exact synthesis interface for matrix-multiplication
circuits with different algebraic requirements at different levels. A
bilinear outer identity is evaluated on matrix-valued blocks, while a
possibly nonbilinear commutative circuit is used only on scalar entries.
If the outer identity has \(r_a\) products and the scalar leaf of order
\(b\) has \(\ell_b\) scheduled products, their explicit composition
computes order \(ab\) with \(r_a\ell_b\) scheduled multiplication gates.
A second interface refines valid matrix-valued instruction programs,
including transpose, while preserving their states, sharing and cost.
The Lean development proves these constructions and closes the explicit
order-16 count 2208 without unproved component premises. A
denominator-cleared realization keeps its intermediates integral and
uses guaranteed exact division by eight only at the outputs. Using the
known commutative leaf bound \(W(b)=b(b^2+2b-1)/2\) yields the general
bound \(r_aW(b)\); perfect squares are one specialization. We give
comparison tests, factor-selection and arbitrary-order fringe
statements, retaining the historical order-81 example and its separate
certificate scope. Scheduled product counts, tensor rank, current
optimality and measured speed remain distinct.

\textbf{Keywords:} exact matrix multiplication; arithmetic circuits;
commutative algorithms; block composition; Lean; certified synthesis.

\subsection{1. Contribution and prior
work}\label{contribution-and-prior-work}

Fixed-size multiplication counts remain useful in exact arithmetic when
coefficient products are expensive. A circuit that exploits
commutativity can save products by multiplying forms that mix entries of
both input matrices. Such a circuit needs a different composition rule
from a usual bilinear tensor decomposition. The scalar calculation and
the matrix-valued block calculation must be separated explicitly.

This article develops the generalization in the author's July 2026
manuscript into a reproducible construction interface. Its central
object is the assembled circuit, with its input maps, scalar products,
output maps and exact schedule cardinality. The general composition
argument is also a short mathematical proof, so the paper can be
reviewed without reading the full Lean implementation.

Block composition and fringe reduction have substantial precedents in
Drevet, Islam and Schost {[}DIS{]}. The commutative leaf formula and its
division-free realization are due to Rosowski {[}R{]}; Waksman's earlier
formulation uses division by two. Rosowski also gives specially
structured nonbilinear algorithms that support recursion through
transpose operations. Thus nonbilinear recursion is possible in
appropriate frameworks. Our condition is narrower: the selected scalar
leaf is not automatically a valid identity on arbitrary noncommuting
blocks.

The 48-product rational outer identity belongs to Dumas, Pernet and
Sedoglavic {[}DPS{]}, and the cited 486-product ternary outer identity
belongs to Perminov {[}P{]}. These are inputs to the synthesis theorem.
We do not claim to have discovered those decompositions, the commutative
leaf, or block multiplication. The contribution of the present revision
is a precise construction and proof interface, its domain and cost
discipline, and the explicit generalized synthesis and comparison
statements. Independent mathematical priority for a numerical portfolio
requires comparisons beyond the older DIS table.

\subsection{2. Model and component
certificates}\label{model-and-component-certificates}

Let \(K\) be a commutative unital ring. Write \(M_b(K)\) for the ring of
\(b\times b\) matrices. Fixed coefficient scales and linear combinations
are recorded separately from multiplication gates whose inputs are
computed from the matrix entries. This is an arithmetic-complexity
model; it does not assign hardware latency to an operation.

An outer scheme of length \(r\) is specified by coefficients
\(\alpha_{\rho ij},\beta_{\rho ij},\gamma_{ij\rho}\in K\). Define

\[
U_\rho(X)=\sum_{i,j<a}\alpha_{\rho ij}X_{ij},\qquad
V_\rho(Y)=\sum_{i,j<a}\beta_{\rho ij}Y_{ij}.
\]

For the required block size, its correctness certificate is the identity

\[
(XY)_{ij}=\sum_{\rho<r}\gamma_{ij\rho}U_\rho(X)V_\rho(Y),
\qquad X,Y\in M_a(M_b(K)).
\tag{1}
\]

Coefficients act centrally on the blocks and the order of the two
factors in every product is retained. Scalar correctness alone is not
substituted for (1). A standard bilinear coefficient/tensor certificate
supplies this block-stability property after the usual
central-coefficient lifting argument; that argument and any particular
input certificate must be checked in their actual coefficient domain.

A scalar leaf has \(\ell\) indexed products. Its left and right input
forms \(L_\lambda(A,B),R_\lambda(A,B)\) may use entries from both input
matrices, and its output recombination is

\[
\operatorname{Leaf}(A,B)_{uv}
=\sum_{\lambda<\ell}w_{uv\lambda}
L_\lambda(A,B)R_\lambda(A,B).
\tag{2}
\]

Its separate certificate is \(\operatorname{Leaf}(A,B)=AB\) for every
\(A,B\in M_b(K)\). The leaf is therefore a circuit for scalar matrix
multiplication; it need not be a bilinear tensor decomposition. The
multiplication schedule is the finite index set \(\{0,\ldots,\ell-1\}\).
Degenerate forms can produce constant or zero gates, so the schedule
size always gives an upper bound on input-dependent products and is not
an optimality assertion.

\subsection{3. Explicit hybrid
composition}\label{explicit-hybrid-composition}

Index an order-\(ab\) matrix by pairs \((i,u)\) with \(i<a,u<b\). The
block map is

\[
\operatorname{blk}(A)_{ij}(u,v)=A_{(i,u),(j,v)}.
\tag{3}
\]

The inverse simply flattens the blocks. Distributing the ordinary
matrix-product sum over the pair index gives

\[
\operatorname{blk}(AB)=\operatorname{blk}(A)\operatorname{blk}(B).
\tag{4}
\]

For each outer product \(\rho\), form the two \(b\times b\) scalar
matrices \(U_\rho(\operatorname{blk}(A))\) and
\(V_\rho(\operatorname{blk}(B))\). Apply the scalar leaf to these two
matrices and recombine its outputs using \(\gamma\). Every actual scalar
product in the result has the explicit index \((\rho,\lambda)\).

\textbf{Theorem 1 (sound composition and scheduled cost).} Suppose the
outer scheme satisfies (1) at block size \(b\), and the leaf (2) is
correct on \(M_b(K)\). The constructed hybrid output is \(AB\). Its
scheduled multiplication-gate set has cardinality \(r\ell\).

\textbf{Proof.} Leaf correctness replaces each constructed leaf output
by the corresponding block product \(U_\rho V_\rho\). The outer
correctness identity then gives the block product of the two blocked
inputs. Equation (4), followed by flattening, identifies this output
with \(AB\). All input and output transformations are fixed linear
combinations. The remaining scheduled products are indexed by the
Cartesian product of the two finite product index sets, whose
cardinality is \(r\ell\). No correctness or gate-count property of the
desired hybrid circuit is used as a hypothesis. \(\square\)

The formal companion implements this construction over arbitrary
commutative rings. Its component hypotheses refer to the outer scheme
and scalar leaf separately. This is a normal compositional proof: it
does not certify an arbitrary external scheme merely because the scheme
declares a product count.

\subsection{4. The scalar commutative
leaf}\label{the-scalar-commutative-leaf}

Rosowski's square-matrix construction supplies the known bound

\[
W(b)=\frac{b(b^2+2b-1)}2,\qquad b\ge2.
\tag{5}
\]

For even \(b\), the mechanism is transparent. Pair adjacent inner
indices \(u=2h,v=2h+1\). For a row \(i\) and a noninitial output column
\(j\), define

\[
\begin{aligned}
P_{ih}&=a_{iu}(b_{u0}+a_{iv}),&
R_{ih}&=a_{iv}(b_{v0}-a_{iu}),\\
Q_{hj}&=b_{vj}(b_{u0}+b_{uj}),&
M_{ihj}&=(a_{iu}+b_{vj})(a_{iv}+b_{u0}+b_{uj}).
\end{aligned}
\tag{6}
\]

The initial output column is \(\sum_h(P_{ih}+R_{ih})\); every other
column is \(\sum_h(M_{ihj}-P_{ih}-Q_{hj})\). Expanding these expressions
in the commutative ring cancels the extraneous same-input products and
leaves the ordinary two summands from each inner-index pair. The \(Q\)
products are shared across rows. The product families have sizes
\(b^2\), \(b(b-1)/2\) and \(b^2(b-1)/2\), whose sum is (5). The
odd-order construction uses Rosowski's separate formulas rather than an
unsupported leftover-pair argument.

At \(b=4\), these four product families contain \(8+8+6+24=46\) gates.
This small leaf is useful as an explicit formal test of the mixed-input
model. Commutativity is essential to its cancellation; applying the same
formulas to arbitrary matrix-valued entries is not justified by the
scalar proof.

\textbf{Corollary 2 (general factor composition).} An admissible outer
scheme of order \(a\) and length \(r_a\), composed with a valid leaf of
length \(W(b)\), supplies an order-\(ab\) multiplication circuit with
scheduled cost

\[
H(a,b)=r_aW(b).
\tag{7}
\]

For perfect-square order \(m^2\), choosing \(a=b=m\) gives \(r_mW(m)\).
For unequal factors, reversing the two chosen sizes can change both the
available coefficient domain and the cost. The synthesis theorem
therefore does not privilege perfect squares.

\subsection{5. Beyond bilinear outer
stages}\label{beyond-bilinear-outer-stages}

The outer stage can be specified by an instruction graph rather than
bilinear coefficient arrays. Use input registers containing the two
blocked matrices' entries or zero. Allow copying, addition, subtraction,
fixed coefficient scaling, block transpose and block multiplication.
Each instruction reads the current register state; chronology preserves
sharing and operation order without expanding the graph into a tree.

\textbf{Theorem 3 (instruction refinement).} Suppose this outer graph
has \(r\) multiplication instructions and, when evaluated on actual
\(b\times b\) matrix blocks, correctly computes the outer matrix
product. Replace each multiplication instruction by an actual
product-vector evaluation of a correct scalar leaf of length \(\ell\),
followed by its fixed linear output recombination. The translated graph
has the same register states at every chronological stage and the same
matrix output. Its scheduled scalar multiplication count is \(r\ell\).

\textbf{Proof.} Initially the register states agree. Copying, linear
operations and transpose preserve this agreement. At a multiplication
instruction, the leaf's scalar matrix-correctness theorem makes its
reconstructed output equal to the original block product. Induction
along the instruction list proves equality of all states and outputs.
Each original multiplication instruction contributes one leaf product
vector of length \(\ell\); all other instructions contribute no new
input-dependent multiplication gate. This gives the stated schedule
count while retaining the original sharing and operation order.
\(\square\)

The companion Lean theorem uses the correctness of the \emph{unrefined}
outer graph and the separate leaf as its component hypotheses, and
derives the refined conclusion. Its transpose instruction is actual
matrix transpose. This broadens the July bilinear-only synthesis
interface to eligible nonbilinear outer programs, including the kind of
involution-sensitive program studied by Rosowski. It does not prove that
every commutative formula is a valid recursive outer program, or claim a
new outer algorithm or improved exponent.

\subsection{6. Comparisons and factor
selection}\label{comparisons-and-factor-selection}

For \(a\ge1,b\ge2\), direct simplification gives

\[
\frac{H(a,b)}{(ab)^3}
=\frac{r_a}{a^3}\left(\frac12+\frac1b-\frac1{2b^2}\right).
\tag{8}
\]

Thus an outer count no worse than \(a^3\) yields a strict saving from
the classical cubic count. The direct commutative construction already
improves on the cubic count, so it is a stronger initial comparator. The
exact test against that construction is

\[
H(a,b)<W(ab)
\iff
r_a<\frac{a(a^2b^2+2ab-1)}{b^2+2b-1}.
\tag{9}
\]

At a perfect square, comparison to using the same outer scheme twice
becomes \(r_mW(m)<r_m^2\) exactly when \(W(m)<r_m\), provided \(r_m>0\).
An independently specified baseline \(B_{ab}\) is beaten precisely when
\(r_aW(b)<B_{ab}\); neither (8) nor (9) proves a current best-known
result.

A catalog of certified outer circuits supplies a finite optimization
problem. Enumerate admissible factor chains \(n=d_1\cdots d_k\), use
block-stable outer circuits at the first \(k-1\) levels, and use the
commutative scalar leaf at the last level. Each chain supplies the upper
bound

\[
\left(\prod_{i<k}r_{d_i}\right)W(d_k).
\tag{10}
\]

The final leaf is the only level allowed to rely on its scalar
commutative cancellation. Taking the smallest admissible recorded cost
selects a witness circuit; it does not establish the globally smallest
possible circuit. Domain compatibility, concrete component validity and
a replayable factor lineage are part of every candidate, rather than
annotations added after a count is selected.

{\def\LTcaptype{none} % do not increment counter
\begin{longtable}[]{@{}
  >{\raggedright\arraybackslash}p{(\linewidth - 6\tabcolsep) * \real{0.1100}}
  >{\raggedright\arraybackslash}p{(\linewidth - 6\tabcolsep) * \real{0.3500}}
  >{\raggedleft\arraybackslash}p{(\linewidth - 6\tabcolsep) * \real{0.1400}}
  >{\raggedright\arraybackslash}p{(\linewidth - 6\tabcolsep) * \real{0.4000}}@{}}
\toprule\noalign{}
\begin{minipage}[b]{\linewidth}\raggedright
Order
\end{minipage} & \begin{minipage}[b]{\linewidth}\raggedright
Selected outer/leaf
\end{minipage} & \begin{minipage}[b]{\linewidth}\raggedleft
Scheduled count
\end{minipage} & \begin{minipage}[b]{\linewidth}\raggedright
Comparison and scope
\end{minipage} \\
\midrule\noalign{}
\endhead
\bottomrule\noalign{}
\endlastfoot
16 & rational outer \(r_4=48\), leaf \(W(4)=46\) & 2208 & 88 below
\(W(16)=2296\); requires the outer dyadic coefficients \\
12 & Strassen outer \(r_2=7\), leaf \(W(6)=141\) & 987 & 15 below
\(W(12)=1002\); integer-coefficient outer \\
81 & cited ternary outer \(r_9=486\), leaf \(W(9)=441\) & 214326 & 57915
below \(W(81)=272241\); cited integer-coefficient certificate \\
\end{longtable}
}

These are valid constructions in their specified models, not assertions
of worldwide optimality. In particular, the familiar 2208-versus-2212
comparison concerns the historical DIS commutative table. The modern
catalogs distinguish fields, commutative and bilinear models, explicit
schemes and formula-only bounds {[}C{]}. A lower count in characteristic
two is not automatically a rational-domain comparator.

A tempting additional factorization is to combine a reported 27-order
outer count 10045 with the order-three commutative leaf count 21, giving
210945 for order 81. The downloaded LRP artifact currently has 10250
active products, with no zero or proportional ordered-product rows
explaining the discrepancy. Its rational coefficients use denominators
dividing 15. Consequently that advertised input has not been promoted to
a verified new corollary here. Recovering the matching certified input
is a concrete further research task, rather than evidence of a completed
improvement.

\subsection{7. An arbitrary-order fringe
construction}\label{an-arbitrary-order-fringe-construction}

\textbf{Theorem 4 (fringe bound).} Let a valid order-\(p\) kernel have
scheduled count \(C_p\le p^3\), with \(p\ge1\). For any \(n\ge1\), let
\(q=\lfloor n/p\rfloor\). There is an order-\(n\) circuit of scheduled
cost

\[
F_p(n)=n^3-q^3(p^3-C_p).
\tag{11}
\]

When \(C_p<p^3\), the construction strictly improves on the cubic count
exactly when \(q>0\), equivalently \(n\ge p\).

\textbf{Proof.} Partition each index set into \(q\) full groups of width
\(p\) and a final group of width \(n-qp\) if it is nonempty. In the
ordinary blocked product, exactly \(q^3\) triples of block indices use
full groups in all three positions. Replace each corresponding cubic
\(p\)-block multiplication by the supplied kernel. Leave the remaining
triples unchanged. The original scalar-product schedule has \(n^3\)
gates and every replacement saves \(p^3-C_p\), proving (11). The output
remains the ordinary block-sum matrix product. \(\square\)

With the order-16 kernel, (11) is \(n^3-1888\lfloor n/16\rfloor^3\). It
applies to prime orders too, without claiming that it beats every other
commutative construction. For instance it gives 3025 at order 17, while
the direct leaf gives 2737. Against that leaf the exact test is

\[
q^3(p^3-C_p)>\frac{n(n-1)^2}{2}.
\tag{12}
\]

Padding provides another candidate, but it must pay for the padded
schedule unless zero simplification is actually performed. Equations
(9), (11) and (12) allow a finite catalog to select exact-factor,
fringe, direct or padded candidates under a common domain and cost
model.

\subsection{8. Integral intermediates and exact output
quotients}\label{integral-intermediates-and-exact-output-quotients}

The recovered 48-product input has integral left and right coefficient
vectors; only its output coefficients require a common denominator
eight. Replace those output coefficients by their integer numerators.
The actual integer coefficient certificate then contracts to eight times
the matrix-product target. Its lifting gives a denominator-cleared block
output \(8XY\) over every commutative ring, with no invertibility
assumption.

Compose that numerator-valued outer stage with the explicit 46-product
leaf. All input linear forms, products and output accumulations have
integer coefficients. The resulting actual order-16 evaluator \(N\)
satisfies

\[
N(A,B)=8AB,
\qquad A,B\in M_{16}(K),
\tag{13}
\]

over every commutative unital ring \(K\). This identity, including the
supplied coefficient and leaf proofs, is now checked in Lean. At integer
inputs, every output entry is divisible by eight and

\[
\operatorname{IntAlg}(A,B)_{ij}=N(A,B)_{ij}\mathbin{/}8=(AB)_{ij},
\qquad A,B\in M_{16}(\mathbb Z).
\tag{14}
\]

Here the quotient is exact integer division, including for negative
outputs. The Lean theorem proves both the actual integer matrix product
and the same 2208-slot multiplication schedule without assuming that two
is a unit in the integers. There are 256 final fixed divisions. They are
explicitly outside the variable-product count and must be included in an
implementation's operation audit.

The same argument applies at order \(4b\) when an actual correct
integral leaf of order \(b\) and length \(\ell\) is provided: the
numerator is eight times the product and the schedule has \(48\ell\)
products. This arbitrary-\(b\) construction and exact integer quotient
are also proved in Lean; leaf validity is its separate normal component
hypothesis, and the desired numerator or final product is derived. With
Rosowski's known leaf this gives \(48W(b)\) integer products followed by
\((4b)^2\) exact output quotients. This is a change in realization and
arithmetic domain, rather than a new product-count or tensor-rank
record. It addresses the July manuscript's restriction that the literal
dyadic reference schedule need not keep integer inputs integral at
intermediate stages. No runtime or bit-complexity improvement is
inferred solely from moving the divisions.

\subsection{9. Verification, numerical behavior and
availability}\label{verification-numerical-behavior-and-availability}

\textbf{Review entry points.} The
\href{https://github.com/alejandrozu/openmath-2026-judging/blob/9e13c89a51a62958ced9ebdf16da31f5e2fb3cdd/personal-matrix-hybrid-2026-10-07/CURRENT_GUIDE.md}{immutable
companion index} binds this manuscript to its actual proof sources.
Start with
\href{https://github.com/alejandrozu/openmath-2026-judging/blob/9e13c89a51a62958ced9ebdf16da31f5e2fb3cdd/personal-matrix-hybrid-2026-10-07/lean/BlockComposition.lean}{block
composition},
\href{https://github.com/alejandrozu/openmath-2026-judging/blob/9e13c89a51a62958ced9ebdf16da31f5e2fb3cdd/personal-matrix-hybrid-2026-10-07/lean/DAGRefinement.lean}{instruction
refinement} and the
\href{https://github.com/alejandrozu/openmath-2026-judging/blob/9e13c89a51a62958ced9ebdf16da31f5e2fb3cdd/personal-matrix-hybrid-2026-10-07/lean/Concrete16.lean}{concrete
order-16 theorem}; then inspect the
\href{https://github.com/alejandrozu/openmath-2026-judging/blob/9e13c89a51a62958ced9ebdf16da31f5e2fb3cdd/personal-matrix-hybrid-2026-10-07/lean/IntegerLateDivision.lean}{integer
output-division theorem} and its
\href{https://github.com/alejandrozu/openmath-2026-judging/blob/9e13c89a51a62958ced9ebdf16da31f5e2fb3cdd/personal-matrix-hybrid-2026-10-07/lean/IntegerGeneralization.lean}{arbitrary-size
generalization}. The
\href{https://github.com/alejandrozu/openmath-2026-judging/blob/9e13c89a51a62958ced9ebdf16da31f5e2fb3cdd/personal-matrix-hybrid-2026-10-07/lean/final_source_manifest.json}{final
11-source/45-declaration manifest} and its checker record exact hashes
and endpoint scope. The
\href{https://github.com/alejandrozu/openmath-2026-judging/blob/9e13c89a51a62958ced9ebdf16da31f5e2fb3cdd/personal-matrix-hybrid-2026-10-07/circuit/README.md}{flat-circuit
replay instructions} offer a separate exact polynomial check. These
links supply reviewable evidence without replacing the readable
component hypotheses and proofs above.

The supplied original reference implementation of the order-16 circuit
is pinned at \texttt{spicylemonade/faster\_16x16}, revision
\texttt{7743ee6848ed876615012c4edb3bc46a1f870afd}. The July audit
independently checked the outer noncommuting polynomial identity, all
4096 output coefficients, the commutative leaf expansion and the
multiplication schedule. It corrected the coefficient-domain wording and
identified an unsafe integer-truncation comparison in the original
verifier. Exact rational equality is the required comparison; truncating
a rational output can conceal an error.

The integer ternary order-nine input is associated with revision
\texttt{3183c54c754b60311edb1947417bc6060bc04db1}; the July record gives
SHA-256
\texttt{6f8ee5d7e89221d2a63c702c0e42444387644c70e08a905bebab172f1033a2b6}.
Its stated 531441 coefficient checks and integer domain belong to that
retained audit. This revision does not relabel historical external
checks as a new Lean replay.

The formal companion's actual checked scope is reported in its
source/endpoint index and compiler receipts. The generic composition
theorem uses outer block-correctness and scalar leaf-correctness as
normal component inputs. Their desired hybrid conclusion and scheduled
cardinality are derived. The particular 48-product input now also has an
actual Lean proof of all 4096 denominator-cleared coefficient equations
and their lifting to every block size over commutative rings in which
two is invertible. The 486-product input retains its separate historical
exact audit; the advertised 10045-product input is not certified here.
The linked records distinguish component validation, generic synthesis
and each concrete assembly.

\textbf{Concrete formal corollary.} The frozen nine-source core closes
the order-16 assembly. For every commutative ring \(K\) in which two is
invertible and all actual \(16\times16\) input matrices, the specified
algorithm equals their matrix product and its scheduled gate count is
2208. This theorem has no external outer-validity, leaf-validity or
coefficient-certificate hypothesis: those component facts are proved
from the supplied data and explicit leaf. Its 35 selected declarations
use only standard kernel axioms. The two additive integer modules close
the numerator, fixed quotient and arbitrary-\(b\) generalization,
bringing the actual selected scope to 11 sources and 45 declarations. No
admitted proof or native-evaluation axiom supplies these results.

An independent generator produces a sparse, fully flattened order-16
circuit with 2208 products in 512 input variables. Its left and right
forms have integer coefficients and its 256 outputs have a common
denominator eight. The separate verifier expands all 256 commutative
quadratic output polynomials and checks them against eight times the
ordinary matrix product. All 4096 expected mixed-input terms match and
every same-input-bank term cancels. The certificate, generator,
independent verifier and exclusive replay receipts are supplied. Their
exact polynomial check took 1.14 seconds in the recorded run; that is
verification time, not multiplication throughput or a formal proof of
the generator itself.

{\def\LTcaptype{none} % do not increment counter
\begin{longtable}[]{@{}
  >{\raggedright\arraybackslash}p{(\linewidth - 2\tabcolsep) * \real{0.2300}}
  >{\raggedright\arraybackslash}p{(\linewidth - 2\tabcolsep) * \real{0.7700}}@{}}
\toprule\noalign{}
\begin{minipage}[b]{\linewidth}\raggedright
Result or artifact
\end{minipage} & \begin{minipage}[b]{\linewidth}\raggedright
Verification scope in this revision
\end{minipage} \\
\midrule\noalign{}
\endhead
\bottomrule\noalign{}
\endlastfoot
Block and instruction-graph composition & General Lean theorem; separate
component hypotheses; actual matrix output and scheduled cost derived \\
Paired-inner commutative leaf & General rectangular even-inner
construction and explicit order-four bridge in Lean \\
48-product outer identity & All 4096 integer coefficient equations and
dyadic lifting to arbitrary block sizes proved in Lean \\
Concrete 2208-product circuit & Unconditional actual order-16 matrix
theorem in Lean, under the commutative-ring/two-invertible domain \\
Integer numerator and final quotients & Actual all-ring \(N=8AB\)
identity and exact integer output-division algorithm proved in Lean;
2208 scheduled products plus 256 fixed quotients \\
General integer order \(4b\) & Actual numerator and quotient
construction in Lean from a separately valid \(\ell\)-product leaf;
\(48\ell\) scheduled products \\
Flat sparse JSON circuit & Independent exact integer polynomial check;
generator itself is not asserted Lean verified \\
Universal odd-order leaf and 486-product input & Published/input
mathematics and retained exact audit; not new full Lean ports \\
Arbitrary-order fringe construction & The readable constructive proof
above; no full encoded fringe generator asserted \\
\end{longtable}
}

The arithmetic count is also distinct from numerical accuracy. In the
order-four leaf, taking \(a_{00}=a_{01}=T\), \(b_{00}=b_{10}=T^{-1}\)
and other entries zero makes the relevant intermediate terms \(1+T^2\)
and \(1-T^2\), whose exact sum is two. A floating-point implementation
can lose the corrections before recombination. The exact proof therefore
supports exact arithmetic without asserting a stable replacement for a
conventional floating-point GEMM kernel.

For applications, additions, coefficient scales and movement must be
counted alongside the saved products. The July literal order-16
schedule, after removal of zero initializations, records 2208 variable
products, 16096 additions/subtractions and 5392 dyadic scales, versus
4096 products and 3840 additions in the classical schedule. Optimizing
the linear network and testing an exact-arithmetic workload are
appropriate next steps; a multiplication count alone does not establish
latency or energy savings.

\textbf{Authorship and acknowledgements.} Alejandro Zarzuelo Urdiales
authored the personal generalized review manuscript. The earlier
concrete result and reference implementation were reported under
Archivara Research Team and \texttt{spicylemonade}; their provenance is
retained. The outer decompositions and commutative leaf retain the
attribution given above. Codex AI assisted the present research
synthesis, drafting and additive formalization; the author must review
the mathematical text, references and contribution statement before
submission. Alejandro also judged and evaluated OpenMath 2026
submissions in the event's Harvard/MIT community context. That judging
role is separate from independent refereeing of this personal paper.

\subsection{References}\label{references}

{[}DIS{]} C.-E. Drevet, M. N. Islam and E. Schost, \emph{Optimization
techniques for small matrix multiplication}, Theoretical Computer
Science 412 (2011), 2219-2236.
\href{https://doi.org/10.1016/j.tcs.2010.12.012}{Publisher article and
correct DOI}.

{[}R{]} A. Rosowski, \emph{Fast commutative matrix algorithms}, Journal
of Symbolic Computation 114 (2023), 302-321.
\href{https://arxiv.org/abs/1904.07683v2}{Earlier full-text version,
arXiv:1904.07683v2}.

{[}DPS{]} J.-G. Dumas, C. Pernet and A. Sedoglavic, \emph{A
non-commutative algorithm for multiplying 4x4 matrices using 48
non-complex multiplications}.
\href{https://arxiv.org/abs/2506.13242v7}{arXiv:2506.13242v7}, 2026
revision.

{[}P{]} A. I. Perminov, \emph{Meta Flip Graph meets Serendipitous
Product: new Fast Matrix Multiplication results}.
\href{https://arxiv.org/abs/2606.02480}{arXiv:2606.02480};
\href{https://github.com/dronperminov/FastMatrixMultiplication}{primary
algorithm repository}.

{[}C{]} B. Chatain Lacelle, \emph{A catalog of fast matrix
multiplication algorithms with exhaustive derivations}.
\href{https://arxiv.org/abs/2606.13408v3}{arXiv:2606.13408v3}. This is
an independent comparator, not the author's repository.

{[}F{]} \href{https://fmm.univ-lille.fr/27x27x27.html}{FMM-Lille
order-27 catalog entry} and its linked coefficient artifact; the
advertised-count/downloaded-schedule discrepancy is recorded in the
supplementary research report.

{[}I{]}
\href{https://github.com/spicylemonade/faster_16x16/blob/7743ee6848ed876615012c4edb3bc46a1f870afd/main.py}{Pinned
original order-16 implementation}.

\end{document}
